% !TeX program = xelatex
\documentclass[UTF8,fontset=fandol,11pt,a4paper]{ctexart}
\usepackage{amsmath,amssymb,bm,geometry,booktabs,longtable,array,xcolor,hyperref}
\geometry{margin=2.25cm}
\hypersetup{colorlinks=true,linkcolor=blue!45!black,urlcolor=blue!45!black}
\setlength{\emergencystretch}{3em}
\linespread{1.12}
\numberwithin{equation}{section}
\newcommand{\Tr}{\operatorname{Tr}}
\newcommand{\diag}{\operatorname{diag}}
\newcommand{\sym}{\operatorname{sym}}
\newcommand{\norm}[1]{\left\lVert #1\right\rVert}
\newcommand{\inner}[2]{\left(#1\,\middle|\,#2\right)}
\newcommand{\code}[1]{\texttt{\detokenize{#1}}}
\title{局域 RI-J 与 COSX 修复\\公式推导、优化条件与数值验收}
\author{C172H346 计算修复记录}
\date{2026 年 11 月 13 日}
\begin{document}
\maketitle
\noindent\textbf{2026-11-18 更新。}
生产修复已改为全局 DF-J（weigend）加 PySCF SGX，
C172H346/cc-pVTZ 收敛到 \(-6717.379184800415\) Ha（Phase 1，16538 s），
加速版 \(-6717.379184800597\) Ha（15315 s，差 \(-1.8\times10^{-10}\) Ha）。
推导与验收见第 \ref{sec:repair} 节。下文为 2026-11-13 的记录，
其中“不改用全局 DF-J”的约束已被该修复路线取代。

\medskip
\noindent\textbf{文档状态：公式推导与修复验收说明。}
2026-11-13 22:40（北京时间）只读复核集群后，最新生产日志是
\path{outputs/c172h346_ccpvtz_5.stdout.log}，修改时间
2026-11-13 20:18:01。生产目录 \code{hf.py} 的 SHA-256 仍为
\texttt{89832978bd7f4feae0c82e601304b25a0ac86d64fc3531470058f0fe331c0f40}，
与出错版本一致。该日志第 1 轮能量为 \(+1.280213786785741\times10^{10}\) Ha，
第 2 轮为 \(+4.412647949115224\times10^{10}\) Ha，文件在第 3 轮开始处结束。
第 1 轮能量与此前 \code{c172h346\_ccpvtz\_4} 打印到同一数值，
说明故障没有换成新的能量公式错误，仍是原局域 J 进入 Fock 后的放大。
能量定义保持
\(E=E_{\mathrm{nuc}}+\Tr(Dh)+\frac12\Tr(DJ)-\frac14\Tr(DK)\)，
不改用全局 DF-J，也不提高辅助基或网格。

同日本地复测：66 项单元测试通过。
C\(_2\)H\(_6\)/STO-3G 在最终 \(35\times302\) 网格收敛，
候选能量 \(-78.3006696922162\) Ha，
PySCF 参考 \(-78.30066969218089\) Ha，
差值 \(3.531\times10^{-11}\) Ha；
独立轨道梯度 RMS \(3.180\times10^{-9}\)，
电子数误差 \(1.066\times10^{-14}\)。
此前 C21/cc-pVTZ 初始密度对照仍有效：
J 最大元素误差约 \(7.49\times10^{-12}\) Ha，
库仑能误差约 \(5.97\times10^{-11}\) Ha。
这些证据不代替 C172 收敛、GPU 对照或近线性标度验收。

\medskip
\noindent\textbf{阅读建议：}
先读第 \ref{sec:roadmap} 节把握计算链；
第 \ref{sec:ri}--\ref{sec:bugs} 节解释 J，
第 \ref{sec:cosx}--\ref{sec:overlap} 节解释 K；
第 \ref{sec:error}、\ref{sec:complexity}、\ref{sec:evidence} 节解释误差与验收。
\begin{abstract}
本文从闭壳层 Hartree--Fock 的密度矩阵出发，推导库仑矩阵 \(J\)、交换矩阵
\(K\)、RI 拟合、层次低秩远场、预条件迭代求解，以及 COSX 的重叠校正。
重点解释原局域实现为什么破坏计算一致性、修复必须保持哪些恒等式，
以及局域和近线性各自需要什么条件。
假设为实数、非相对论、闭壳层、有限分子高斯基组。
数学目标不等于实现已经通过验收；原始故障、设计条件与新实现状态分别记录。
\end{abstract}
\tableofcontents
\newpage

\section{先把计算链连接起来}
\label{sec:roadmap}
每轮 SCF 的依赖关系是
\begin{equation}
D\longrightarrow\{J(D),K(D)\}\longrightarrow F=h+J-\tfrac12K
\longrightarrow C_{\mathrm{occ}}\longrightarrow D_{\mathrm{new}}.
\label{eq:flow}
\end{equation}
错误 \(J\) 会改变 Fock 和占据轨道，再通过新密度影响下一轮。
因此增加迭代次数不能修复一个已经错误的积分映射。

RI-J 必须保持
\begin{equation}
\boxed{\gamma=B^Td,\qquad Mc=\gamma,\qquad j=Bc.}
\label{eq:roadJ}
\end{equation}
这里 \(d,j\) 按第 \ref{sec:packing} 节打包。
优化的是算子存储和作用方式，而非改变这三条关系。

本次实际参考 SGX 默认采用双边对称重叠拟合：
\begin{equation}
Q=S_{\mathrm{num}}^{-1}S,\quad P=\tfrac12(Q+I),\qquad
\boxed{K=\sym\!\left(P^T\mathcal R(PD)\right).}
\label{eq:roadK}
\end{equation}
\(\mathcal R\) 是第 \ref{sec:cosx} 节的网格交换映射。
两个 \(P\) 的位置都要保留，不能只做原子对角块修正。

正确性和性能分别验收：最终
\(|E-E_{\mathrm{ref}}|\leq10^{-6}\) Ha；
指定烷烃序列的 J、K 实测标度指数分别不超过 1.5。
能量通过不自动证明近线性，计时通过也不证明物理结果正确。

\section{符号、指标与单位}
\begin{longtable}{p{0.19\textwidth}p{0.71\textwidth}}
\toprule 符号&含义\\\midrule\endhead
\(\mu,\nu,\lambda,\sigma\)&AO 指标，范围 \(1,\ldots,n\)。\\
\(P,Q\) 作为下标&辅助基指标，范围 \(1,\ldots,m\)。\\
\(i,g\)&占据轨道指标、网格点指标。\\
\(\phi_\mu,\chi_P\)&实数 AO、实数辅助基函数。\\
\(D,S,h,F\)&自旋求和密度、重叠、单电子 Hamiltonian、Fock 矩阵。\\
\(M,B,\gamma,c\)&辅助库仑度量、打包三中心积分、拟合右端、拟合系数。\\
\(X,W,S_{\mathrm{num}}\)&网格 AO 值、权重对角矩阵、数值重叠。\\
\(Q,P\) 作为矩阵&完整重叠投影、对称 SGX 投影；见式 \eqref{eq:roadK}。\\
\(\sym(A)\)&\((A+A^T)/2\)。\\
\(\norm{A}_F,\norm{A}_2\)&Frobenius 范数、谱范数。\\
\bottomrule
\end{longtable}
使用原子单位：长度为 Bohr，能量为 Hartree（Ha）。
Angstrom 输入在分子构建时统一转换。
通常 AO 矩阵约定下 \(D,S,C\) 无量纲，\(h,F,J,K\) 数值单位为 Ha。
辅助函数可以采用不同归一化，但 \(B,M,c\) 必须一起一致变换。

\section{从 RHF 密度得到 Fock 与总能量}
闭壳层每个占据空间轨道含两个电子：
\begin{equation}
D_{\mu\nu}=2\sum_{i=1}^{n_{\mathrm{occ}}}C_{\mu i}C_{\nu i},
\qquad C_{\mathrm{occ}}^TSC_{\mathrm{occ}}=I.
\end{equation}
因此轨道更新后的密度满足
\begin{equation}
\Tr(DS)=2n_{\mathrm{occ}}=N_e,\qquad DSD=2D.
\label{eq:count}
\end{equation}
电子数不是 \(\Tr(D)\)。原子叠加初猜未必严格幂等，
应与占据轨道产生的密度分开检查。

定义双电子积分
\begin{equation}
(\mu\nu|\lambda\sigma)=
\iint\frac{\phi_\mu(\bm r)\phi_\nu(\bm r)
\phi_\lambda(\bm r')\phi_\sigma(\bm r')}
{|\bm r-\bm r'|}\,d\bm r\,d\bm r'.
\end{equation}
库仑与交换矩阵是
\begin{align}
J_{\mu\nu}(D)&=\sum_{\lambda\sigma}D_{\lambda\sigma}
(\mu\nu|\lambda\sigma),\label{eq:Jdef}\\
K_{\mu\nu}(D)&=\sum_{\lambda\sigma}D_{\lambda\sigma}
(\mu\lambda|\nu\sigma).\label{eq:Kdef}
\end{align}
自旋求和约定对应
\begin{align}
F&=h+J-\tfrac12K,\\
E&=E_{\mathrm{nuc}}+\Tr(Dh)+\tfrac12\Tr(DJ)-\tfrac14\Tr(DK)
\label{eq:energy}\\
 &=E_{\mathrm{nuc}}+\tfrac12\Tr[D(h+F)].
\end{align}
库仑项的 \(1/2\) 消除电子对重复计数，交换项额外的 \(1/2\)
来自闭壳层自旋求和。若改为单自旋密度，必须同时改变这些系数。
能量在轨道正交约束下取极值，得到
\begin{equation}
FC=SC\varepsilon.
\end{equation}
错误 Fock 也可能产生严格满足式 \eqref{eq:count} 的轨道；
电子数正确只检查归一化，不能替代 J/K 正确性检查。

\section{RI-J 的变分推导}
\label{sec:ri}
\subsection{在库仑内积中拟合密度}
电子密度及辅助展开为
\begin{equation}
\rho(\bm r)=\sum_{\mu\nu}D_{\mu\nu}\phi_\mu(\bm r)\phi_\nu(\bm r),
\qquad \widetilde\rho(\bm r)=\sum_Pc_P\chi_P(\bm r).
\end{equation}
定义
\begin{equation}
\inner{u}{v}=\iint\frac{u(\bm r)v(\bm r')}{|\bm r-\bm r'|}
\,d\bm r\,d\bm r',
\quad
M_{PQ}=\inner{\chi_P}{\chi_Q},
\quad
b_{\mu\nu,P}=\inner{\phi_\mu\phi_\nu}{\chi_P}.
\end{equation}
辅助函数在库仑内积下线性无关时 \(M\) 对称正定。
真实线性相关和数值病态须显式处理，不能无记录地添加对角偏移。

最小化拟合误差
\begin{align}
\mathcal L(c)&=\inner{\rho-\widetilde\rho}{\rho-\widetilde\rho}\\
&=\inner{\rho}{\rho}-2c^T\gamma+c^TMc,\\
\gamma_P&=\sum_{\mu\nu}D_{\mu\nu}b_{\mu\nu,P}.
\end{align}
对系数求导：
\begin{equation}
\frac{\partial\mathcal L}{\partial c}=2Mc-2\gamma=0
\quad\Longrightarrow\quad Mc=\gamma.
\label{eq:normal}
\end{equation}
求解线性方程即可，不需要显式形成 \(M^{-1}\)。

\subsection{AO 对打包与二重权重}
\label{sec:packing}
令 \(p=(\mu,\nu)\)、\(\mu\geq\nu\)，定义
\begin{equation}
w_p=\begin{cases}1,&\mu=\nu,\\2,&\mu>\nu,\end{cases}
\qquad
d_p=\sqrt{w_p}D_{\mu\nu},\quad
B_{pP}=\sqrt{w_p}\,b_{\mu\nu,P}.
\label{eq:pack}
\end{equation}
于是
\begin{equation}
(B^Td)_P=\sum_{\mu\geq\nu}w_pb_{\mu\nu,P}D_{\mu\nu}
=\sum_{\mu\nu}b_{\mu\nu,P}D_{\mu\nu}=\gamma_P.
\end{equation}
令 \(j_p=\sqrt{w_p}J_{\mu\nu}\)，则
\begin{equation}
j=Bc,\qquad J_{\mu\nu}=j_p/\sqrt{w_p},\qquad
J_{\nu\mu}=J_{\mu\nu}.
\label{eq:unpack}
\end{equation}
因此
\begin{equation}
E_J^{\mathrm{RI}}=\tfrac12d^Tj
=\tfrac12\gamma^TM^{-1}\gamma=\tfrac12c^TMc.
\label{eq:rijenergy}
\end{equation}

\textbf{与实际实现的映射。}
当前 J 模块采用等价的原始积分打包：
\begin{equation}
B^{\mathrm{code}}_{pP}=b_{\mu\nu,P},\quad
d^{\mathrm{code}}_p=w_pD_{\mu\nu},\quad
j^{\mathrm{code}}_p=J_{\mu\nu}.
\end{equation}
它仍满足 \(\gamma=(B^{\mathrm{code}})^Td^{\mathrm{code}}\)、
\(j^{\mathrm{code}}=B^{\mathrm{code}}c\)、
\(E_J=\tfrac12(d^{\mathrm{code}})^Tj^{\mathrm{code}}\)。
若 \(W_p=\diag(w_p)\)，
\begin{equation}
B=W_p^{1/2}B^{\mathrm{code}},\quad
d=W_p^{-1/2}d^{\mathrm{code}},\quad
j=W_p^{1/2}j^{\mathrm{code}}.
\end{equation}
正文的平方根权重使普通欧氏转置直接对应密度空间的伴随；
实现无需真的额外乘平方根，但所有能量和互易性检查须使用一致权重。

\subsection{对称性、互易性和能量正性}
三步合并：
\begin{equation}
j=\mathcal A d,\qquad \mathcal A=BM^{-1}B^T.
\label{eq:operator}
\end{equation}
若 \(M\) 正定，
\begin{equation}
\mathcal A^T=\mathcal A,\qquad
d^T\mathcal A d=(B^Td)^TM^{-1}(B^Td)\geq0,
\end{equation}
并有
\begin{equation}
d_1^Tj(d_2)=d_2^Tj(d_1).
\label{eq:reciprocity}
\end{equation}
互易性比 AO 矩阵 \(J=J^T\) 更强。
事后转置平均不能修复密度到 J 的错误映射。

同一固定密度下
\begin{equation}
E_J^{\mathrm{exact}}-E_J^{\mathrm{RI}}
=\tfrac12\inner{\rho-\widetilde\rho}{\rho-\widetilde\rho}\geq0.
\end{equation}
本次实现误差对照相同辅助基组的 RI-J；
RI 相对四中心积分的模型误差另行报告。

\section{原局域 J 为什么会出错}
\label{sec:bugs}
\subsection{密度元素为零不代表输出 J 为零}
式 \eqref{eq:Jdef} 表明
\begin{equation}
D_{\mu\nu}=0\quad\not\Longrightarrow\quad J_{\mu\nu}=0,
\end{equation}
因为 \(J_{\mu\nu}\) 接收所有 \(D_{\lambda\sigma}\) 的作用。
密度可筛选 \(\gamma=B^Td\) 的收缩工作；
输出 \(j=Bc\) 必须按拟合系数与输出误差独立处理。

H$_2$/STO-3G、键长 \(1\) \AA、weigend、\(D=I\) 的已复现结果是
\begin{align}
J_{\mathrm{ref}}=
\begin{pmatrix}1.252553646591&0.618671094989\\
0.618671094989&1.252553646591\end{pmatrix},\\
J_{\mathrm{old}}=
\begin{pmatrix}1.252553646591&0\\0&1.252553646591\end{pmatrix}.
\end{align}
该算例只有一个区域，全部辅助壳层保留，已排除空间截断。
但因 \(D_{12}=0\)，漏掉的非对角项不直接贡献该密度的 \(\Tr(DJ)\)；
这个反例不能单独解释 C21 库仑能放大的全部量值。

\subsection{高斯局域不等于库仑势没有远场}
若局域分布总电荷 \(q=\int u\,d\bm r\)，其远场通常为
\begin{equation}
V_u(\bm R)=q/|\bm R|+O(|\bm R|^{-2}).
\end{equation}
高斯函数值衰减快，但其库仑势具有长程尾部。
辅助函数中心不在邻域内，不足以作为删除全部三中心库仑贡献的理由。
即使单极为零，仍可能有偶极及更高阶贡献。

\subsection{局域求逆与 owned 拼接不保证一致}
旧实现对带 halo 的区域求逆，再只输出该区域所属辅助函数。
其输入投影、求逆区域和输出投影一般不互为伴随，
总体映射通常不能写成同一个 \(BM^{-1}B^T\)。
局部度量各自正定并不保证拼接后满足式 \eqref{eq:reciprocity}。
修复保留局域区域用于精确近场和预条件，
用同一个压缩算子保留跨区域耦合。
各个缺陷对异常量值的贡献仍需消融实验，不能把一条条件当成完整定量解释。

\section{近场精确、远场层次低秩}
\label{sec:hierarchy}
对辅助函数和 AO 乘积建立空间树。
几何及积分误差控制决定可忽略的 AO 乘积，输出空间不随密度零元素改变。
近场块使用精确积分，充分分离的远场块尝试
\begin{equation}
A_{ts}\approx U_{ts}V_{ts}^T,\qquad
U_{ts}\in\mathbb R^{n_t\times r},\quad
V_{ts}\in\mathbb R^{n_s\times r}.
\end{equation}
库仑核在分离区域的平滑性提供低秩机会，但并不保证任意块都低秩。
Gaussian 非紧支撑，空间半径本身不是能量误差定理。

初始参数为叶节点最多 128 个自由度、远场分离系数 2、
最大块秩 64、容差 \(10^{-12}\)。
秩不足时细分，精度不足时收紧至 \(10^{-13}\)、\(10^{-14}\)。
近场可以流式重新计算并使用有界缓存；数学定义不要求长期存下全部近场块。

设 ACA 残差 \(R^{(k)}=A-A^{(k)}\)，选非零主元 \((i_k,j_k)\)：
\begin{align}
u_k&=R^{(k)}(:,j_k),\\
v_k^T&=R^{(k)}(i_k,:)/R^{(k)}(i_k,j_k),\\
A^{(k+1)}&=A^{(k)}+u_kv_k^T.
\end{align}
每步查询必要行列的精确积分，避免预先形成整个远场块。
需要主元保护、确定性抽样检查；必要时可重压缩。
当前实现另取独立的保留行列检查拟合，若检查失败，将新行列加入主元搜索，
继续增秩；达到最大秩仍不足时细分块。只在参与拟合的行列上检查残差
会漏掉未采样的方向，本次 C21 修复过程中已实际遇到这一问题。
ACA 抽样残差不是严格谱误差上界，不能替代独立全量参考验收。

正向和转置必须共用因子：
\begin{equation}
y_t\mathrel{+}=U_{ts}(V_{ts}^Tx_s),\qquad
z_s\mathrel{+}=V_{ts}(U_{ts}^Tx_t).
\label{eq:adjoint}
\end{equation}
度量 \(M\) 的对称块成对应用同一份因子。
对称不等于正定，仍需线性求解和能量检查。

\section{缩放、两层预条件与 PCG}
\label{sec:pcg}
令
\begin{equation}
L=\diag(\sqrt{M_{11}},\ldots,\sqrt{M_{mm}}),\quad
A=L^{-1}ML^{-1},\quad b=L^{-1}\gamma,\quad y=Lc.
\end{equation}
原问题转为 \(Ay=b\)、\(c=L^{-1}y\)。
变量缩放不改变模型，也不能消除真实线性相关。

令 \(R_s\) 提取区域，\(W_s\) 为重叠权重，\(A_s=R_sAR_s^T\)。
对称加性 Schwarz 为
\begin{equation}
H_{\mathrm{local}}=\sum_sR_s^TW_sA_s^{-1}W_sR_s.
\end{equation}
左右对称使用权重，所有自由度必须被覆盖。
当前实现使用经典对称加性 Schwarz，即 \(W_s=I\)。
重叠次数归一化并非必须；即使另一组权重仍保持正定，
也可能恶化预条件效果。权重只影响求解效率，不应改变被求解的全局方程。
加入粗空间 \(Z\)：
\begin{equation}
A_c=Z^TAZ,\qquad H=H_{\mathrm{local}}+ZA_c^{-1}Z^T.
\label{eq:precond}
\end{equation}
粗层处理长波误差。
一般理论方案需要粗层大小或其递归作用保持受控；
当前首版采用有上限的聚合粗空间，不能仅凭“两层”二字保证链长无关的迭代数，
必须实测 PCG 次数与体系大小的关系。

取 \(r_k=b-Ay_k,z_k=Hr_k\)，PCG 为
\begin{align}
p_0&=z_0,&
\alpha_k&=\frac{r_k^Tz_k}{p_k^TAp_k},\\
y_{k+1}&=y_k+\alpha_kp_k,&
r_{k+1}&=r_k-\alpha_kAp_k,\\
z_{k+1}&=Hr_{k+1},&
\beta_k&=\frac{r_{k+1}^Tz_{k+1}}{r_k^Tz_k},\\
p_{k+1}&=z_{k+1}+\beta_kp_k.
\end{align}
一次求解内算子、预条件与筛选固定；改变压缩精度后重新启动。
除递推残差外，重算真实残差 \(b-Ay_k\)。
非正曲率、非有限值、停滞和超限必须报告，不能用静默对角修正隐藏。

\section{残差小不等于能量误差已经合格}
\label{sec:error}
暂时假设 \(B,M\) 精确。令
\begin{equation}
c_\star=M^{-1}\gamma,\quad r=\gamma-M\widehat c,\quad
e=c_\star-\widehat c=M^{-1}r.
\end{equation}
二次目标 \(\Phi(c)=\tfrac12c^TMc-c^T\gamma\) 满足
\begin{equation}
\Phi(\widehat c)-\Phi(c_\star)
=\tfrac12e^TMe=\tfrac12r^TM^{-1}r.
\label{eq:resquad}
\end{equation}
稳健拟合能量为
\begin{equation}
E_{\mathrm{rob}}(\widehat c)
=\gamma^T\widehat c-\tfrac12\widehat c^TM\widehat c
=E_\star-\tfrac12r^TM^{-1}r.
\end{equation}
但直接由 \(j=B\widehat c\) 组装的普通 RI 能量是
\begin{equation}
E_{\mathrm{assembled}}=\tfrac12\gamma^T\widehat c,\qquad
E_\star-E_{\mathrm{assembled}}=\tfrac12\gamma^TM^{-1}r.
\label{eq:reslinear}
\end{equation}
不能把式 \eqref{eq:resquad} 的二次残差误差直接套到式 \eqref{eq:reslinear}。
Cauchy--Schwarz 给出
\begin{equation}
|E_\star-E_{\mathrm{assembled}}|
\leq\tfrac12\sqrt{\gamma^TM^{-1}\gamma}\sqrt{r^TM^{-1}r}.
\end{equation}
若知道 \(\lambda_{\min}(M)>0\)，可用
\(r^TM^{-1}r\leq\norm{r}_2^2/\lambda_{\min}(M)\)。
相对残差不能直接读作 Hartree 误差。

层次表示还引入 \(\widetilde B=B+\Delta B\)、\(\widetilde M=M+\Delta M\)。
近似方程解得精确不代表算子准确。
先令 \(\Delta B=0\)，若
\(\eta=\norm{M^{-1/2}\Delta M M^{-1/2}}_2<1\)，则
\begin{equation}
\left|\tfrac12\gamma^T(\widetilde M^{-1}-M^{-1})\gamma\right|
\leq\frac{\eta}{2(1-\eta)}\norm{M^{-1/2}\gamma}_2^2.
\end{equation}
还需计入 \(\Delta B\) 对右端与输出的影响。
固定密度下
\begin{equation}
|\delta E_J|=\tfrac12|\Tr(D\,\delta J)|
\leq\tfrac12\norm{D}_F\norm{\delta J}_F.
\end{equation}
因此独立矩阵与能量参考不可被 ACA 或 PCG 内部指标替代。

\section{从四中心交换积分推导 COSX}
\label{sec:cosx}
定义网格位置的解析势积分
\begin{equation}
A_{\nu\sigma}(\bm r_g)
=\int\frac{\phi_\nu(\bm r')\phi_\sigma(\bm r')}
{|\bm r_g-\bm r'|}\,d\bm r'.
\label{eq:potential}
\end{equation}
带网格参数的三指标 \(A_{\nu\sigma}(\bm r_g)\)
与前面的缩放辅助度量矩阵不同。
重排式 \eqref{eq:Kdef}：
\begin{align}
K_{\mu\nu}
&=\int d\bm r\ \phi_\mu(\bm r)
 \sum_{\lambda\sigma}\phi_\lambda(\bm r)
 D_{\lambda\sigma}A_{\nu\sigma}(\bm r)\\
&\approx\sum_gw_gX_{g\mu}\sum_\sigma A_{\nu\sigma}(\bm r_g)
 \sum_\lambda X_{g\lambda}D_{\lambda\sigma}.
\end{align}
对任意矩阵 \(Y\) 定义
\begin{align}
f_{g\sigma}[Y]&=\sum_\lambda X_{g\lambda}Y_{\lambda\sigma},\\
G_{g\nu}[Y]&=\sum_\sigma A_{\nu\sigma}(\bm r_g)f_{g\sigma}[Y],\\
\mathcal R(Y)&=X^TWG[Y],\qquad W=\diag(w_g).
\label{eq:raw}
\end{align}
不做重叠校正时 \(K=\sym(\mathcal R(D))\)。
全量 \(A_{g\nu\sigma}\) 需 \(O(N_gn^2)\) 存储，
应按网格和壳层对现场计算并立即收缩。
分批减少峰值内存，总工作量仍取决于访问了多少壳层对。

\subsection{将积分与收缩融合执行}
当前实现复用 PySCF/libcint 的 C 层
\code{SGXnr_direct_k_drv}，融合执行
\(G_{g\nu}=\sum_\sigma A_{\nu\sigma}(\bm r_g)f_{g\sigma}\)，
避免在 Python 中形成网格乘两个 AO 指标的完整张量。
本程序仍负责投影密度、局域支持和筛选、原始
\(R=X^TWG\) 积累以及 \(P^T\) 输出校正；
生产路径不调用高层 \code{get_k_only} 或 SCF 参考求解器。
C 层积分收缩在 CPU 执行，GPU 只计入实际的 CuPy 运算阶段。
这个改动改变执行方式，式 \eqref{eq:raw} 的数学定义保持不变。
参考与候选共享底层积分原语，因此独立对照覆盖投影、组装和 SCF，
不宣称独立验证该共享 C 内核。显式张量分块模式仍用于复核融合原语。

\section{SGX 重叠拟合：默认需要双边对称校正}
\label{sec:overlap}
\subsection{完整投影与非对称选项}
数值重叠和解析重叠分别为
\begin{equation}
S_{\mathrm{num}}=X^TWX,\qquad
S_{\mu\nu}=\int\phi_\mu\phi_\nu\,d\bm r.
\end{equation}
求
\begin{equation}
S_{\mathrm{num}}Q=S,\qquad Q=S_{\mathrm{num}}^{-1}S,
\label{eq:Q}
\end{equation}
则 \(\overline X=XQ\) 满足 \(X^TW\overline X=S\)。
这是离散交叉重叠条件，不代表所有交换积分精确。
PySCF 的非对称拟合选项为
\begin{equation}
K_{\mathrm{unsym}}=\sym(\mathcal R(QD)).
\end{equation}

\subsection{实际默认与一阶误差消除}
实施检查确认，实际 SGX 对象默认开启对称拟合；
不能把内部数据类构造函数的默认值当成实际配置。
默认使用
\begin{equation}
P=\tfrac12(Q+I),\qquad
K_{\mathrm{sym}}=\sym(P^T\mathcal R(PD)).
\label{eq:symK}
\end{equation}
它相当于数值积分一侧的两个 AO 因子都使用 \(XP\)。
令 \(\Delta=Q-I\)。因 \(S_{\mathrm{num}}Q=S\)，有
\begin{align}
P^TS_{\mathrm{num}}P
&=S_{\mathrm{num}}
+\tfrac12\Delta^TS_{\mathrm{num}}
+\tfrac12S_{\mathrm{num}}\Delta
+\tfrac14\Delta^TS_{\mathrm{num}}\Delta\\
&=S+\tfrac14\Delta^TS_{\mathrm{num}}\Delta.
\end{align}
对称校正将重叠误差消到一阶，剩余项为二阶。
它不同于 \(Q\) 的精确交叉重叠关系。

这是实施过程中对初始方案的明确修正：
只用 \(S_{\mathrm{num}}^{-1}SD\) 对应非对称路径，
不能直接与默认对称 SGX 做精密对照。
两种模式都须使候选与参考选项一致；本次默认采用式 \eqref{eq:symK}。

\subsection{不形成稠密投影矩阵的等价求解}
输入投影可写为
\begin{equation}
S_{\mathrm{num}}Z_D=SD,\qquad
Y=PD=\tfrac12(Z_D+D).
\label{eq:inputsolve}
\end{equation}
令 \(R=\mathcal R(Y)\)，输出投影为
\begin{equation}
S_{\mathrm{num}}Z_R=R,\qquad
P^TR=\tfrac12(SZ_R+R),\qquad
K=\sym\!\left[\tfrac12(SZ_R+R)\right].
\label{eq:outputsolve}
\end{equation}
用到 \(Q^T=SS_{\mathrm{num}}^{-1}\)，要求两种重叠矩阵对称。
这允许逐列或逐块稀疏求解。

\(PD\) 一般不对称，不能在输入时先对称化。
原始矩阵方向必须是 \(R=X^TWG\)；
若保存的是 \(R^T=G^TWX\)，对应投影应转到右侧，不能仍左乘 \(P^T\)。
没有输出投影时，转置后对称化等价；
加入投影后，中间方向会影响结果。

旧实现只修正每个原子的 \(K_{aa}\)，缺少跨原子重叠、逆作用和密度耦合。
局域化应通过稀疏支持与完整方程残差实现，
不能先删掉跨原子项再假定与全局投影等价。

\section{局域支持、残差扩域与增量 K}
\label{sec:screen}
对任意投影方程 \(S_{\mathrm{num}}Z=B_{\mathrm{rhs}}\)，
近似解 \(\widehat Z\) 满足
\begin{equation}
R_Z=B_{\mathrm{rhs}}-S_{\mathrm{num}}\widehat Z,\qquad
Z-\widehat Z=S_{\mathrm{num}}^{-1}R_Z.
\end{equation}
残差必须包含区域外的行；只检查局域子问题会漏掉边界误差。
扩域仍无法达到目标时应失败或报告标度退化，不能静默放宽精度。
两次投影和全部筛选误差都应传播到
\begin{equation}
|\delta E_K|=\tfrac14|\Tr(D\,\delta K)|
\leq\tfrac14\norm{D}_F\norm{\delta K}_F.
\end{equation}

原始网格贡献包含
\begin{equation}
w_gX_{g\mu}A_{\nu\sigma}(\bm r_g)f_{g\sigma}.
\end{equation}
可按绝对贡献上界筛选：
\begin{equation}
|\Delta R_{\mu\nu}|
\leq\sum_{\text{被忽略 }(g,\sigma)}
|w_g|\,|X_{g\mu}|\,|A_{\nu\sigma}(\bm r_g)|\,|f_{g\sigma}|.
\end{equation}
这约束原始 \(R\)，最终还需考虑输出投影。
网格权重可能有负值，上界使用 \(|w_g|\)。
应区分网格附近 AO 支持、密度耦合后的 \(f=XY\) 支持、
以及解析势和输出支持。三者使用同一裸半径会漏项。
交换局域性依赖 AO、密度衰减与贡献上界共同作用，
不能对金属或任意随机稠密密度承诺相同标度。

固定网格、投影、积分与筛选时
\begin{equation}
K(D_n)=K(D_{n-1})+K(\Delta D_n).
\end{equation}
若掩码依赖输入，一般
\begin{equation}
K_{\mathcal S(D_n)}(D_n)\ne
K_{\mathcal S(D_{n-1})}(D_{n-1})
+K_{\mathcal S(\Delta D_n)}(\Delta D_n).
\end{equation}
增量 K 应由完整密度决定支持，掩码变化时完整重建或补偿。
每五轮及收敛前完整重建是控制策略，不替代完整 K 自身精度验证。
最终细网格改变 \(X,W,S_{\mathrm{num}},P\)，
需要清除不兼容缓存与 DIIS，重新构建并继续收敛。
只在旧密度上重算一次能量不能称为细网格自洽。

\section{近线性优化的条件与计时口径}
\label{sec:complexity}
设体系规模 \(N\)，固定基组下 \(n,m,N_g=O(N)\)。
若固定精度时有效 AO 乘积数为 \(O(N)\)、远场平均秩为 \(r\)，
典型层次存储与应用成本为
\begin{equation}
{\rm storage}=O(rN\log N),\qquad
{\rm matvec}=O(rN\log N).
\end{equation}
若 PCG 次数为 \(q\)，J 成本必须累计为
\begin{equation}
T_J=T_{B^Td}+qT_{M\,\mathrm{apply}}
+qT_{\mathrm{precond}}+T_{Bc}.
\end{equation}
秩、邻域、预条件成本与迭代次数都要受控。
仅看到局域积分循环不能推出整个求解近线性。

K 还要求每批 AO、密度耦合、积分对与投影求解工作受控。
计时包含筛选、积分、搬运和 GPU 同步。
稠密 J/K 接口展开与存储为 \(O(n^2)\)，
稠密广义本征求解主要成本为 \(O(n^3)\)；
它们单独报告，完整 SCF 不宣称近线性。

相邻实际完成规模的指数是
\begin{equation}
p=\log(T_2/T_1)/\log(n_2/n_1).
\end{equation}
C21/C43/C86/C172 使用同设备同精度，预热一次、测三次取中位数。
初始化与稳态时间分开，每次重置线性求解初值，
避免相同右端已收敛解制造零迭代加速。
报告秩、PCG 次数、峰值内存和稠密展开。
指定序列上的实测结论不是任意分子的复杂度定理。

\section{原实现证据与修复验收分开}
\label{sec:evidence}
C21H44/cc-pVTZ/weigend 同一正常初猜的独立诊断为
\begin{align}
\max|D_0|&=1.9249278484,\\
\lambda_{\min}(S)&=1.63808007\times10^{-4},\\
E_{J,\mathrm{old}}&=230710.83169559\ {\rm Ha},\\
E_{J,\mathrm{ref}}&=1839.34748314\ {\rm Ha},\\
|E_{K,\mathrm{old}}-E_{K,\mathrm{ref}}|
&\simeq7.40\times10^{-4}\ {\rm Ha}.
\end{align}
J 误差约 \(2.2887\times10^5\) Ha，主导该对照。
错误初始 Fock 第一次对角化后
\begin{equation}
\max|D_1|=2337.73048640,\qquad
\Tr(D_1S)=169.99999999999.
\end{equation}
这显示反馈放大，并说明电子数正确仍不足以认证结果。
C172 原日志 \code{c172h346\_ccpvtz\_4} 第 13 轮约为
\(+1.35901636\times10^{11}\) Ha，未收敛；
该任务在北京时间 2026-11-13 11:30:50 已停止。
更新的 \code{c172h346\_ccpvtz\_5}（2026-11-13 20:18:01）仍由同一
\texttt{89832978} 生产源码生成，第 1 轮能量与上一份日志相同，为
\(+1.280213786785741\times10^{10}\) Ha，第 2 轮为
\(+4.412647949115224\times10^{10}\) Ha，并在第 3 轮开始处中断。
两次运行的首轮能量一致，说明巨大正值来自同一局域 J 故障，
而不是能量组装符号或一次偶然的浮点异常。
这些均为原实现数据，不能作为新算法通过的证据。

验收依次为：
\begin{enumerate}
\item 小体系完整 \(B,M\) 对照：伴随、互易、线性、正性和压缩误差。
\item 固定密度 J/K 最大元素误差各 \(\leq10^{-8}\) Ha；
库仑能和交换能误差各 \(\leq10^{-7}\) Ha。
\item 双方完整 SCF 收敛，能量差 \(\leq10^{-6}\) Ha；
能量变化和轨道梯度 RMS 的 \(10^{-8}\) 门槛不放宽。
\item 更新后密度电子数误差 \(\leq10^{-6}\)，检查对称性和幂等性。
\item 最大档 J/K 时间指数各 \(\leq1.5\)，C172 必须实际完成。
\end{enumerate}
覆盖平移、旋转、远离二聚体、零块密度和原空间分支边界。
随机密度测试算子正确性，绝缘烷烃测试性能，不共享无条件标度承诺。

\section{公式如何对应修改流程}
\begin{longtable}{p{0.23\textwidth}p{0.31\textwidth}p{0.35\textwidth}}
\toprule 步骤&数学关系&应检查的输出\\\midrule\endhead
积分与打包&式 \eqref{eq:pack}--\eqref{eq:unpack}&完整小体系积分、二重权重。\\
层次远场&式 \eqref{eq:adjoint}&共用正向/转置因子、块秩。\\
拟合求解&式 \eqref{eq:normal}、\eqref{eq:precond}&真实残差、正曲率、迭代数。\\
J 与能量&式 \eqref{eq:operator}、\eqref{eq:rijenergy}&独立 J、互易性、库仑能。\\
COSX&式 \eqref{eq:raw}&原始矩阵方向、网格收缩。\\
双边投影&式 \eqref{eq:inputsolve}--\eqref{eq:outputsolve}&全局残差、跨原子项、同模式 SGX。\\
增量与网格&固定映射线性、缓存一致性&完整重建差值、细网格再收敛。\\
性能归因&各阶段成本累计&时间、内存、秩、次数、标度。\\
\bottomrule
\end{longtable}
实际代码与日志以最终源码哈希为准。
公式定义和设计条件不替代实现验证，尚未满足的条件必须在状态页写明。

\section{修复与加速：完整公式推导（2026-11-18 更新）}
\label{sec:repair}
本节对应生产代码 \code{hf.py} 中的
\code{scf_energy_terms}、\code{check_physical_energy}、
\code{build_df_j}、\code{build_sgx_k}、\code{DenseLinearAlgebra}、
\code{job_k_memory_mb}、\code{_sgx_block_size} 与网格切换逻辑。
前文第 \ref{sec:ri}--\ref{sec:bugs} 节给出的是局域 RI-J 的设计目标；
本节记录实际采用的修复路线：全局 DF-J 加 PySCF SGX（COSX），
以及在其上不改变能量的加速。

\subsection{闭壳层能量分解与物理下界}
\label{sec:repair-energy}
对占据系数 \(C\equiv C_{\mathrm{occ}}\in\mathbb R^{n\times n_{\mathrm{occ}}}\)，
\begin{equation}
D=2CC^T,\qquad C^TSC=I,\qquad
E=E_{\mathrm{nuc}}+\underbrace{\Tr(Dh)}_{E_1}
+\underbrace{\tfrac12\Tr(DJ)}_{E_J}
\underbrace{-\tfrac14\Tr(DK)}_{E_K}.
\label{eq:rep-energy}
\end{equation}
代码逐项计算这四个量，并另算 \(\Tr(DS)\)。

\paragraph{\(E_J>0\)：库仑度量正定。}
记 \(\rho(\bm r)=\sum_{\mu\nu}D_{\mu\nu}\phi_\mu(\bm r)\phi_\nu(\bm r)\)，则
\begin{equation}
\Tr(DJ)=\sum_{\mu\nu\lambda\sigma}D_{\mu\nu}D_{\lambda\sigma}(\mu\nu|\lambda\sigma)
=\inner{\rho}{\rho}
=\int\frac{4\pi}{k^2}\,|\hat\rho(\bm k)|^2\,\frac{d\bm k}{(2\pi)^3}>0,
\label{eq:rep-ej}
\end{equation}
因为 \(1/r\) 的 Fourier 变换 \(4\pi/k^2\) 处处为正；
只要 \(\rho\not\equiv0\) 就严格为正。
DF 近似下，令拟合右端
\(\gamma_P=\sum_{\mu\nu}(P|\mu\nu)D_{\mu\nu}\)
（任务记录中记作 \(d_P\)），\(M_{PQ}=(P|Q)\) 对称正定，则
\begin{equation}
E_J^{\mathrm{DF}}=\tfrac12\Tr(DJ^{\mathrm{DF}})
=\tfrac12\sum_{\mu\nu}D_{\mu\nu}\sum_P(\mu\nu|P)c_P
=\tfrac12\gamma^TM^{-1}\gamma
=\tfrac12\norm{L^{-1}\gamma}_2^2\geq0,
\label{eq:rep-ejdf}
\end{equation}
其中 \(M=LL^T\) 为 Cholesky 分解。
正性不依赖密度是否幂等，也不依赖辅助基是否完备；
它只依赖“同一个正定 \(M\) 同时出现在拟合和回代中”。
这正是第 \ref{sec:repair-root} 小节局域实现失去的性质。

\paragraph{\(E_K<0\)：交换收缩半正定。}
把 \(D=2CC^T\) 代入式 \eqref{eq:Kdef}：
\begin{equation}
\Tr(DK)=\sum D_{\mu\nu}D_{\lambda\sigma}(\mu\lambda|\nu\sigma)
=4\sum_{i,j=1}^{n_{\mathrm{occ}}}
\inner{\varphi_i\varphi_j}{\varphi_i\varphi_j}\geq0,
\qquad \varphi_i=\sum_\mu C_{\mu i}\phi_\mu .
\label{eq:rep-ek}
\end{equation}
每一项都是“交换电荷” \(\varphi_i\varphi_j\) 的库仑自能，
由式 \eqref{eq:rep-ej} 同理非负，且 \(i=j\) 项严格为正。
因此 \(E_K=-\sum_{ij}\inner{\varphi_i\varphi_j}{\varphi_i\varphi_j}<0\)。
若 \(D\) 只是半正定（如初猜），用谱分解
\(D=\sum_kd_kv_kv_k^T\)（\(d_k\geq0\)）可得同样结论：
\(\Tr(DK)=\sum_{kl}d_kd_l\inner{v_kv_l}{v_kv_l}\geq0\)。
同时 \(K\) 作为 \(D\) 的线性映射在 Frobenius 内积下半正定，
SGX 近似只要保持对称，就不应使 \(E_K\) 变号。

\paragraph{电子数。}
\(\Tr(DS)=2\Tr(C^TSC)=2n_{\mathrm{occ}}=N_e\)。
C\(_{172}\)H\(_{346}\) 有 \(N_e=6\times172+346=1378\)，
cc-pVTZ 下 \(n=30\times172+14\times346=10004\)。

\paragraph{MINAO 初猜为何不检查电子数。}
MINAO 先在原子最小基中放入球平均占据 \(o_k\)，
再把原子轨道 \(\psi_k\) 以最小二乘投影到目标基：
\(\tilde\psi_k=\Pi\psi_k\)，\(\Pi=S^{-1}S_{12}\)（不再归一化），
\(D_0=\sum_ko_k\tilde c_k\tilde c_k^T\)。投影是压缩的，
\(\norm{\Pi\psi_k}^2\leq\norm{\psi_k}^2=1\)，所以
\begin{equation}
\Tr(D_0S)=\sum_ko_k\norm{\Pi\psi_k}^2\leq\sum_ko_k=N_e .
\end{equation}
C172 实测 \(\Tr(D_0S)=1377.9487\)，比 1378 少 \(0.0513\)，
平均每个原子约 \(10^{-4}\)。\(D_0\) 也不满足 \(D_0SD_0=2D_0\)。
因此初始 Fock 只检查 \(E_J>0\)、\(E_K<0\)、\(E<0\)；
第一次对角化后 \(D=2CC^T\) 精确满足电子数，
此后每轮都检查 \(|\Tr(DS)-N_e|\leq10^{-6}\)。

\subsection{根因：局域 RI-J 失去正定度量}
\label{sec:repair-root}
v3 GPU 路线的自定义局域 RI-J 有两处偏离式 \eqref{eq:rep-ejdf}：
\begin{enumerate}
\item 输出端复用了密度筛选：\(D_{\mu\nu}\) 小的位置直接令 \(J_{\mu\nu}=0\)，
等价于在输出加掩码 \(\Pi_{\mathrm{out}}\)；
\item 每个区域 \(R\)（含 halo）各自求逆 \(M_R^{-1}\)，
只把“owned”辅助函数的系数拼回全局。
\end{enumerate}
写成算子：令 \(R_R\) 为限制到区域 \(R\) 的选择矩阵，
\(E_R\) 为只保留 owned 行的延拓，
\begin{equation}
c_{\mathrm{loc}}=G\gamma,\qquad
G=\sum_RE_RM_R^{-1}R_R,\qquad
J_{\mathrm{loc}}=\Pi_{\mathrm{out}}\,B\,G\,B^Td .
\label{eq:rep-gloc}
\end{equation}
一般 \(E_R\neq R_R^T\)，所以 \(G\neq G^T\)，且 \(G\neq M^{-1}\)：
\(c_{\mathrm{loc}}\neq(P|Q)^{-1}(Q|D)\)。由它算出的库仑能
\begin{equation}
E_J^{\mathrm{loc}}=\tfrac12\gamma^TG\gamma
=\tfrac12\gamma^T\sym(G)\gamma
\end{equation}
只由 \(\sym(G)\) 决定。各块 \(M_R^{-1}\) 正定并不保证 \(\sym(G)\) 正定，
其特征值可以为负，也可以远大于 \(\lambda_{\max}(M^{-1})\)。
所以 \(E_J\geq0\) 的保证不再成立，量值也没有上界；
\(\Pi_{\mathrm{out}}\) 还会进一步破坏互易性式 \eqref{eq:reciprocity}。

\paragraph{误差界：变分（robust）拟合与局域拟合。}
令 \(c^\ast=M^{-1}\gamma\)，任取系数 \(c=c^\ast+\delta c\)。
Dunlap 的 robust 泛函
\begin{equation}
\widetilde E_J[c]=\gamma^Tc-\tfrac12c^TMc
=E_J^{\mathrm{DF}}-\tfrac12\delta c^TM\delta c
\label{eq:rep-robust}
\end{equation}
对 \(\delta c\) 只有二阶误差，且符号确定（总是低估）。
同样有 \(E_J^{\mathrm{exact}}-E_J^{\mathrm{DF}}
=\frac12\inner{\rho-\widetilde\rho}{\rho-\widetilde\rho}\geq0\)。
局域实现的 \(J=Bc_{\mathrm{loc}}\) 对应的是非 robust 形式
\(\frac12\gamma^Tc_{\mathrm{loc}}\)，误差是一阶的：
\begin{equation}
E_J^{\mathrm{loc}}-E_J^{\mathrm{DF}}=\tfrac12\gamma^T\delta c,
\qquad \delta c=(G-M^{-1})\gamma .
\end{equation}
在 \(M\) 度量下用 Cauchy--Schwarz：
\(|\gamma^T\delta c|\leq\norm{\gamma}_{M^{-1}}\norm{\delta c}_M\)，
\(\norm{\delta c}_M\leq\norm{M^{1/2}GM^{1/2}-I}_2\norm{\gamma}_{M^{-1}}\)，
且 \(\norm{\gamma}_{M^{-1}}^2=2E_J^{\mathrm{DF}}\)，于是
\begin{equation}
\frac{|E_J^{\mathrm{loc}}-E_J^{\mathrm{DF}}|}{E_J^{\mathrm{DF}}}
\leq\norm{M^{1/2}GM^{1/2}-I}_2 .
\label{eq:rep-bound}
\end{equation}
\(G=M^{-1}\) 时右端为零；局域求逆只在 \(G\) 接近 \(M^{-1}\)
（即谱意义下逼近）时才可控，“每块都正定”不足以保证这一点。

\paragraph{C21 与 C172 的数值。}
C21/cc-pVTZ 初始密度下，局域 J 给出 \(E_J=230711\) Ha，
参考 DF-J 为 \(1839\) Ha，比值约 125。
由式 \eqref{eq:rep-bound}，
\(\norm{M^{1/2}GM^{1/2}-I}_2\geq(230711-1839)/1839\approx124\)，
说明局域度量算子远离 \(M^{-1}\)，而不是普通截断误差。
错误 \(J\) 进入 \(F=h+J-\frac12K\) 后，占据轨道按扭曲的谱排序；
一阶响应 \(\delta D\) 的分母是轨道能差 \(\varepsilon_a-\varepsilon_i\)，
\(\delta J\) 与能隙同量级以上时线性响应不再成立，
新密度把电荷堆到错误区域，下一轮 \(E_J\) 继续放大。
C172 第 1 轮能量因此为 \(+1.280\times10^{10}\) Ha，
第 2 轮 \(+4.41\times10^{10}\) Ha。
加入 \(E_J>0\)、\(E_K<0\)、\(E<0\) 守卫后，此类状态会在第一轮被拒绝。

\subsection{修复：全局 DF-J 与 SGX/COSX}
\label{sec:repair-fix}
\paragraph{DF-J（weigend 辅助基）。}
\begin{align}
\gamma_P&=\sum_{\mu\nu}(P|\mu\nu)D_{\mu\nu}, \label{eq:rep-dfj1}\\
LL^T&=M,\quad L\,y=\gamma,\quad L^Tc=y
\quad(\Rightarrow c=M^{-1}\gamma), \label{eq:rep-dfj2}\\
J_{\mu\nu}&=\sum_P(\mu\nu|P)\,c_P . \label{eq:rep-dfj3}
\end{align}
三中心积分按 AO 对 \(\mu\geq\nu\) 打包，\(D\) 非对角元乘 2（第 \ref{sec:packing} 小节）。
式 \eqref{eq:rep-dfj1} 与 \eqref{eq:rep-dfj3} 各 \(O(N_{\mathrm{aux}}n^2)\)，
Cholesky 为 \(O(N_{\mathrm{aux}}^3)\) 但只做一次，三角回代 \(O(N_{\mathrm{aux}}^2)\)。
同一个 \(M\) 出现在拟合和回代中，式 \eqref{eq:rep-ejdf} 自动成立。
代码 \code{build_df_j} 调用 PySCF \code{df_jk.get_j}（\code{hermi=1}，
\code{direct_scf_tol=1e-13}）。

\paragraph{SGX/COSX 交换。}
在 \(\bm r\) 上做数值积分、在 \(\bm r'\) 上解析积分：
\begin{align}
K_{\mu\nu}&=\sum_{\lambda\sigma}D_{\lambda\sigma}
\iint\frac{\phi_\mu(\bm r)\phi_\lambda(\bm r)\phi_\nu(\bm r')\phi_\sigma(\bm r')}{|\bm r-\bm r'|}
\approx\sum_gw_g\phi_\mu(\bm r_g)\sum_{\lambda\sigma}\phi_\lambda(\bm r_g)D_{\lambda\sigma}A_{\nu\sigma}(g),\\
A_{\nu\sigma}(g)&=\int\frac{\phi_\nu(\bm r')\phi_\sigma(\bm r')}{|\bm r'-\bm r_g|}\,d\bm r' .
\end{align}
令 \(X_{g\mu}=\phi_\mu(\bm r_g)\)，三步为
\begin{equation}
F_{g\sigma}=w_g\sum_\lambda X_{g\lambda}D_{\lambda\sigma},\qquad
G_{g\nu}=\sum_\sigma A_{\nu\sigma}(g)F_{g\sigma},\qquad
K_{\mu\nu}=\sum_gX_{g\mu}G_{g\nu}.
\label{eq:rep-sgx}
\end{equation}
第一、三步是 GEMM，各 \(O(N_gn^2)\)；
第二步需要每个格点的 ESP 积分 \(A(g)\)，形式上也是 \(O(N_gn^2)\)，
由于高斯乘积 \(\phi_\nu\phi_\sigma\) 的筛选，有效对数随 \(n\) 线性增长，
总成本接近 \(O(N_gn)\)，而 \(N_g\propto N_{\mathrm{atom}}\)。
数值 \(K\) 不严格对称，结果取 \(\sym(K)\)，
并用重叠拟合（第 \ref{sec:overlap} 节）消除网格一阶误差。
\paragraph{两级网格。}
先用粗网格 \(30\times110\)（径向 \(\times\) 角向）收敛，
再在细网格 \(35\times302\) 上从粗网格密度出发、用完整 J/K 重新收敛，
最后再做一次完整 J/K 重建。最终能量只属于细网格。

\paragraph{增量 K。}
\(K\) 对 \(D\) 线性，故
\(K(D_{k})=K(D_{k-1})+K(\Delta D_k)\)，\(\Delta D_k=D_k-D_{k-1}\)。
增量步只对 \(\Delta D\) 做 SGX，筛选阈值 \(10^{-10}\)，
筛选所用的壳对仍由完整 \(D_k\) 决定，使增量与完整 K 用同一组壳对。
累积的筛选误差每 \(k_{\mathrm{full}}=5\) 轮清零一次（阈值 \(10^{-13}\) 的完整 K）。
若增量步已满足收敛判据，先重建完整 K 再判停，
避免“\(\Delta D\) 小所以看起来已收敛”的假象。

\paragraph{物理守卫（每轮执行）。}
代码函数为 \code{check_physical_energy}：
\begin{enumerate}
\item \(E,D,F\) 全部有限；
\item \(E_J=\frac12\Tr(DJ)>0\)；
\item \(E_K=-\frac14\Tr(DK)<0\)；
\item \(E<0\)（束缚分子）；
\item 第一次对角化之后 \(|\Tr(DS)-N_e|\leq10^{-6}\)；
\item 增量 K 判停前完整重建并重新检查；最终完整 J/K 重建后再检查一次。
\end{enumerate}
任一条失败即抛出异常并打印各能量分量，不进入下一轮。

\subsection{加速：不改变能量的四项措施}
\label{sec:repair-speed}
\paragraph{(1) K 内存预算决定 SGX 块大小。}
SGX 每个格点同时保存 4 个长度为 \(n\) 的双精度数组，
即每格点 \(b=8\times4\times n\) 字节。
\code{job_k_memory_mb} 给出固定预算 \(M_K\)（\code{k_memory_gib}）
与进程上限 \(M_{\max}=0.75\,M_{\mathrm{job}}\)；
\code{_sgx_block_size} 取
\begin{align}
M_{\mathrm{avail}}&=\min\!\bigl(M_K,\;M_{\max}-M_{\mathrm{RSS}}-M_J\bigr),
\qquad M_J=(n_{\mathrm{thr}}+1)\cdot8n^2,\\
\mathrm{blksize}&=\min\!\Bigl(N_g,\;
B_0\Bigl\lfloor\frac{\max(B_0,\lfloor M_{\mathrm{avail}}/b\rfloor)}{B_0}\Bigr\rfloor\Bigr),
\end{align}
\(B_0\) 为 PySCF 的 \code{SGX_BLKSIZE} 对齐单位。
未设 \(M_K\) 时退回 PySCF 原规则 \(M_{\mathrm{avail}}=M_{\mathrm{max\_memory}}-M_{\mathrm{RSS}}\)：
DIIS 历史增长后 RSS 上升，块被压到最小值，GEMM 效率下降。
C172 时 \(b\approx320\) kB，每 GB 约 3100 个格点。
块大小只改变循环划分，不改变求和结果（至多浮点舍入顺序）。

\paragraph{(2) 粗网格提前切换。}
粗网格阶段改为 \(|\Delta E|<10^{-5}\) Ha 即切换，不再收敛到 \(10^{-8}\)。
细网格 SCF 的不动点 \(D^\ast_{\mathrm{fine}}\) 与初始密度无关（同一收敛域内），
细网格仍按 \(\tau=10^{-8}\) 的能量与梯度判据收敛并做最终完整重建。
SCF 能量对密度变分稳定：
\(E(D^\ast+\delta)-E(D^\ast)=O(\norm{\delta}^2)\)，
所以起点差别只影响迭代数，最终能量差受 \(\tau\) 控制。
C172 实测两种协议最终能量相差 \(-1.8\times10^{-10}\) Ha。

\paragraph{(3) 稠密线性代数：密度一次 GEMM，DIIS 误差两次 GEMM。}
密度 \(D=2C_{\mathrm{occ}}C_{\mathrm{occ}}^T\) 只需一次
\(n\times n_{\mathrm{occ}}\) 乘 \(n_{\mathrm{occ}}\times n\) 的 GEMM，
约 \(2n^2n_{\mathrm{occ}}\) flop。
DIIS 误差 \(e=FDS-SDF\)：由于 \(F,D,S\) 均对称，
\begin{equation}
(FDS)^T=S^TD^TF^T=SDF
\quad\Longrightarrow\quad
H=(FD)S,\qquad e=H-H^T,
\end{equation}
只需两次 \(n^3\) GEMM，原写法需四次；转置是零成本的视图。
新误差向量只与已存向量做内积，DIIS Gram 矩阵每轮增加一行而不是重算。

\paragraph{(4) 为什么不把 SGX 的 GEMM 搬到 GPU。}
SGX 计时中 ESP 积分 \(A(g)\) 占 \(K\) 时间的 84--87\%，GEMM 只占其余部分。
Blackwell GPU 双精度峰值约 0.98 TFLOP/s，与 CPU 的 0.93 TFLOP/s 相当，
GEMM 加速比 \(s\approx1.05\)。按 Amdahl 定律
\begin{equation}
\frac{t_{\mathrm{old}}}{t_{\mathrm{new}}}
=\frac{1}{f+(1-f)/s}
\approx\frac{1}{0.84+0.16/1.05}\approx1.008 ,
\end{equation}
上限不到 1\%，还要加上主机--设备拷贝，所以保留 CPU 路线。

\subsection{验收与标度}
\label{sec:repair-validation}
\paragraph{链长外推。}
烷烃每增加一个 CH\(_2\) 单元，能量近似增加常数
\(\Delta_{\mathrm{CH_2}}=-39.04779646\) Ha，
\begin{equation}
E(\mathrm C_{172})\approx E(\mathrm C_{21})+151\,\Delta_{\mathrm{CH_2}}
=-821.1618495437+151\times(-39.04779646)
=-6717.37912\ \mathrm{Ha}.
\end{equation}
实算值与外推值相差 \(-7.0\times10^{-5}\) Ha（约 0.04 kcal/mol），
说明 C172 结果与短链一致，而失败的 \(+1.28\times10^{10}\) Ha 差了十个数量级以上。

\begin{center}
\small
\begin{tabular}{p{3.1cm}p{3.6cm}p{5.2cm}r}
\toprule
体系/协议 & 能量 / Ha & 对照 & 墙钟 / s\\
\midrule
C21，修复后 & \(-821.1618495437\) & 参考 \(-821.1618494204\)，差 \(-1.23\times10^{-7}\) & --\\
C43，Phase 1 \(\to\) 加速 & -- & -- & 2191.6 \(\to\) 1959.9\\
C86，Phase 1 \(\to\) 加速 & -- & -- & 5813.6 \(\to\) 5466.0\\
C172，Phase 1 & \(-6717.379184800415\) & 外推 \(-6717.3791\) & 16538\\
C172，加速 & \(-6717.379184800597\) & 与 Phase 1 差 \(-1.8\times10^{-10}\) & 15315\\
\bottomrule
\end{tabular}
\end{center}
加速版在 C43、C86、C172 上分别节省 10.6\%、6.0\%、7.4\% 墙钟时间，
能量差在 SCF 收敛阈值以下。

\paragraph{标度指数。}
设 \(t=aN^p\)，取两个规模 \(N_1,N_2\)：
\begin{equation}
\frac{t_2}{t_1}=\Bigl(\frac{N_2}{N_1}\Bigr)^p
\quad\Longrightarrow\quad
p=\frac{\ln(t_2/t_1)}{\ln(N_2/N_1)} .
\end{equation}
以碳数 \(N_1=43,N_2=86\)（基函数数 2522、5016，比值 1.99，结果几乎相同）：
Phase 1 为 \(\ln(5813.6/2191.6)/\ln2=1.41\)，
加速版为 \(\ln(5466.0/1959.9)/\ln2=1.48\)，总体约 1.45。
分项看，PySCF 路径 J 约 \(N^{2.1}\)、K 约 \(N^{1.5}\)；
自研层次 RI 路径 J 预处理约 \(N^{2.3}\)、J 构建约 \(N^{1.7}\)、K 约 \(N^{2.1}\)。
因此层次路径目前没有达到近线性目标，生产仍用 PySCF 路径。
两点拟合只给局部指数；常数项和预处理的 \(O(N^3)\) 部分
（如 \(M\) 的 Cholesky、\(n^3\) 的对角化）在更大体系中会逐渐占主导，
需要三个以上规模做对数线性回归才能判断渐近指数。

\appendix
\section{Gaussian 乘积如何给出距离相关筛选}
两个 primitive 的指数部分满足 Gaussian 乘积定理：
\begin{align}
p&=\alpha+\beta,\qquad \kappa=\frac{\alpha\beta}{p},\qquad
\bm P=\frac{\alpha\bm A+\beta\bm B}{p},\\
e^{-\alpha|\bm r-\bm A|^2}e^{-\beta|\bm r-\bm B|^2}
&=e^{-\kappa|\bm A-\bm B|^2}e^{-p|\bm r-\bm P|^2}.
\end{align}
令 \(t=|\bm r-\bm P|\)，角动量多项式的绝对值可用
\begin{equation}
(t+|\bm P-\bm A|)^{l_a}(t+|\bm P-\bm B|)^{l_b}
=\sum_n b_nt^n,\qquad b_n\geq0
\end{equation}
控制。对各项分别使用
\begin{align}
\int_{\mathbb R^3}t^n e^{-pt^2}\,d\bm r
&=2\pi\,\Gamma\!\left(\frac{n+3}{2}\right)p^{-(n+3)/2},\\
\sup_{t\geq0}t^ne^{-pt^2}
&=\left(\frac{n}{2ep}\right)^{n/2}\quad(n>0),
\end{align}
其中 \(n=0\) 时第二式取 1。
乘上 primitive 绝对系数、保守角向归一化界和
\(e^{-\kappa|\bm A-\bm B|^2}\)，再对所有 primitive 对求和，
可得到下节使用的 \(L_1\) 和 \(L_\infty\) 上界。
距离相关的 Gaussian 乘积衰减控制的是 AO 乘积本身；
它不允许丢掉已保留乘积分布的长程库仑势。

\section{交换筛选中解析势的一个统一上界}
\label{sec:potentialbound}
为避免把裸距离当作库仑误差界，可对
\(u(\bm r)=\phi_\mu(\bm r)\phi_\nu(\bm r)\) 使用
\begin{equation}
L_1=\int |u(\bm r)|\,d\bm r,\qquad
L_\infty=\sup_{\bm r}|u(\bm r)|.
\end{equation}
以任意网格点 \(\bm r_g\) 为中心，用半径 \(a>0\) 将积分分成两部分：
\begin{align}
|A_{\mu\nu}(\bm r_g)|
&\leq
\int_{|\bm r-\bm r_g|\leq a}\frac{|u(\bm r)|}{|\bm r-\bm r_g|}
\,d\bm r
+\int_{|\bm r-\bm r_g|>a}\frac{|u(\bm r)|}{|\bm r-\bm r_g|}
\,d\bm r\\
&\leq 2\pi a^2L_\infty+\frac{L_1}{a}.
\end{align}
令导数为零得 \(a^3=L_1/(4\pi L_\infty)\)，因此
\begin{equation}
\boxed{|A_{\mu\nu}(\bm r_g)|
\leq \tfrac32(4\pi L_\infty L_1^2)^{1/3}.}
\label{eq:uniformV}
\end{equation}
当 \(u=0\) 时上界取零，不执行除零运算。
AO 的 Cauchy--Schwarz 给出
\begin{equation}
L_1\leq\sqrt{S_{\mu\mu}S_{\nu\nu}},\qquad
L_\infty\leq U_\mu U_\nu,\quad
U_\mu\geq\sup_{\bm r}|\phi_\mu(\bm r)|.
\end{equation}
实现可由收缩 Gaussian 各 primitive 的绝对系数、径向最大值及保守角向上界
构造 \(U_\mu\)。这些界需要覆盖所用的笛卡尔或球谐归一化，不能忽略收缩系数。

记式 \eqref{eq:uniformV} 的上界为 \(V_{\mu\nu}^{\max}\)。
若在一个网格批次内忽略 \(f\) 的第 \(\nu\) 列，
其对原始交换矩阵最大元素的贡献可界为
\begin{equation}
\max_{\lambda\mu}|\Delta R_{\lambda\mu}|
\leq
\left(\max_\mu V_{\mu\nu}^{\max}\right)
\sum_{g\ {\rm in\ batch}}|w_g|
\left(\max_\lambda|X_{g\lambda}|\right)|f_{g\nu}|.
\end{equation}
在批次和被忽略列之间累计预算，才可将逐次筛选转换为全局原始矩阵误差控制。
当前实现为该类 \(f\) 列筛选分配 \(10^{-12}\) 的原始矩阵预算。
该界只控制这一种筛选误差：
它不自动覆盖 AO 空间尾部、投影求解、输出投影放大或最终 SCF 能量。
因此独立 SGX 与总能量验收仍然必要。

\section*{复核资料}
\addcontentsline{toc}{section}{复核资料}
\begin{enumerate}
\raggedright
\item 原诊断目录：\path{outputs/c172_energy_diagnosis_20261113/}，
含 \code{diagnosis_zh.md}、\code{local_kernel_numeric_checks.json}、
\code{local_j_density_screen.json}。
\item 原代码快照：
\path{outputs/c172_repair_20261113/baseline/hf.py} 及 SHA-256 清单。
\item PySCF 2.14.0：
\code{pyscf.sgx.sgx_jk.SGXData._build_ovlp_fit}、
\code{get_k_only}、SGX 实例的 \code{_symm_ovlp_fit} 配置、
\code{pyscf.df.df_jk.get_j}。验收需保存实际版本和模式。
\item Q. Sun 等，Recent developments in the PySCF program package，
\emph{J. Chem. Phys.} \textbf{153}, 024109 (2020)。
本文具体运算顺序以实际检查源码为准。
\end{enumerate}
\end{document}
